\documentclass[10pt,a4paper]{article}
\usepackage[utf8]{inputenc}
\usepackage[T1]{fontenc}
\usepackage[french]{babel}
\usepackage[top=2.1cm, bottom=1.8cm, left=1.8cm, right=1.8cm]{geometry}
\usepackage{enumitem}
\usepackage{fancyhdr}
\setlength{\parindent}{0pt}
\setlength{\parskip}{0.35em}
\pagestyle{fancy}
\fancyhf{}
\lhead{Victor Henri Willy Eder\\Liam Van Camp}
\rhead{27 mars 2026\\EPFL Microtechnique}
\renewcommand{\headrulewidth}{0pt}
\setlength{\headheight}{28pt}
\setlength{\headsep}{18pt}

\begin{document}
\begin{center}
    {\huge \textbf{Rapport d'activit\'e -- Rendu 1}}
\end{center}
\vspace{0.5cm}
{\large \textbf{Part I : Méthode de travail}}

\begin{list}{}{\leftmargin=0.5cm \rightmargin=0cm}
\item
Notre méthode de travail a débuté par une phase de planification
commune, durant laquelle nous avons réparti les différents modules
du projet entre les deux membres du groupe. Nous avons également mis
en place un dépôt Git partagé ainsi que le \textbf{\texttt{Makefile}} (initial), afin
de structurer le développement et faciliter la compilation.

Dans un second temps, chacun a développé ses modules de manière
autonome, en les testant régulièrement et en maintenant le dépôt Git
à jour.

Une fois ces modules implémentés, nous avons collaboré sur
l’intégration finale, en particulier sur les modules \textbf{\texttt{game}}
et \textbf{\texttt{project}}, incluant le fichier \textbf{\texttt{main}}. En complément
du temps de travail en classe, nous nous sommes également réunis
environ une fois par semaine afin d’échanger sur le code, corriger
les éventuelles erreurs et assurer l’harmonisation de nos
contributions.


\end{list}

{\large \textbf{Part II : Répartition du travail}}

\begin{list}{}{\leftmargin=0.5cm \rightmargin=0cm}
\item
    {\normalsize \textbf{Liam Van Camp}}

    \begin{list}{}{\leftmargin=0.5cm \rightmargin=0cm}
    \item
    D\'eveloppement et tests du module \textbf{\texttt{tools}},
    de la classe \textbf{\texttt{Paddle}}
    et de la fonction \textbf{\texttt{main}},
    ainsi que des fonctions li\'ees
    \`a la lecture des fichiers d'entr\'ee.
    \end{list}

    {\normalsize \textbf{Victor Henri Willy Eder}}

    \begin{list}{}{\leftmargin=0.5cm \rightmargin=0cm}
    \item
    D\'eveloppement et tests des classes
    \textbf{\texttt{Ball}} et \textbf{\texttt{Brick}},
    ainsi que des fonctions li\'ees
    aux tests d'intersection
    et d'inclusion.
    \end{list}

    {\normalsize \textbf{Travail effectu\'e ensemble}}

    \begin{list}{}{\leftmargin=0.5cm \rightmargin=0cm}
    \item
    Travail commun sur le module \textbf{\texttt{Game}}.
    Dans ce cadre, Liam s'est principalement charg\'e
    de la lecture du fichier et de la validation des donn\'ees,
    tandis que Victor s'est principalement charg\'e
    des tests d'intersection et d'inclusion.
    En parall\`ele, nous avons aussi effectu\'e
    des ajustements communs du \textbf{\texttt{Makefile}}
    et validation finale sur les tests.
    \end{list}
\end{list}
\end{document}
